\subsection{Discussion}

The present study provides a mechanistic framework linking catalysis, sequence background, and evolutionary history in GPX6. The residue-49 swap agrees with the expected chemical advantage of selenium in GPX active sites, since Sec lowers the proton-transfer barrier relative to Cys in both homologous backgrounds.\cite{Reich2016,prabhakar_elucidation_2005,cardey_selenocysteine_2007,rees_ancient_2024,Kanavos2026SecSelectivity} The important point is that this advantage is not the same in the two scaffolds. In the human background, replacing Sec with Cys raises the barrier by 1.43~kcal/mol. In the mouse background, the reciprocal Cys-to-Sec replacement lowers the barrier by 2.84~kcal/mol. Thus, the Sec-containing state is energetically favored for forming the reactive chalcogenolate in the GPX6 active site, but the size of that benefit is set by the surrounding active-site and near-active-site environment.\cite{Warshel2006,Kamerlin2010,VillaWarshel2001}

The mutational landscapes further show that this catalytic network is epistatic. Mouse-to-human substitutions tend to lower or preserve the activation barrier once the catalytic residue has been swapped, whereas human-to-mouse substitutions destabilize the catalytic-residue-swapped background during the early steps. The same sequence differences therefore produce different energetic outcomes depending on direction, consistent with background-dependent path constraints in protein evolution.\cite{Weinreich2006,starr2016epistasis,storz_compensatory_2018,Kaltenbach2015} The residue-level analysis reinforces this interpretation because substitutions can change magnitude or even sign as the background accumulates additional changes. These inversions are not explained simply by the number of local hydrogen bonds, pointing instead to a cooperative catalytic architecture in which multiple residues jointly determine transition-state stabilization. However, the greedy trajectories should be interpreted as parsimonious barrier-minimizing hypotheses, not as inferred historical substitution orders. They ignore population-genetic effects, lineage-specific constraints outside the 20-residue set, and the possibility that a globally plausible evolutionary route may include locally nonoptimal energetic steps.

Projecting the EVB-derived trajectories onto the GPX6 phylogeny helps interpret this asymmetry in evolutionary terms. Rees et al. inferred five independent Sec losses in mammalian GPX6 and found that Sec-to-Cys branches had elevated evolutionary rates and convergent changes, especially in the GPX catalytic domain.\cite{rees_ancient_2024} They therefore interpreted the Cys-containing lineages as products of convergent adaptation shaped by pleiotropy and epistasis rather than as simple restorations of ancestral Sec-dependent catalysis. Our results provide an energetic counterpart to that observation. The extant human wild-type enzyme appears to occupy a Sec-dependent regime, whereas the mouse enzyme follows a Cys-adapted regime that can regain low proton-transfer barriers when Sec is reintroduced in a favorable background. The triangular projection tempers the expectation of strong hysteresis because the forward and reverse paths overlap in reduced sequence space rather than forming completely separated routes. However, the EVB rankings still reveal epistatic constraints and a partial consensus in the order by which substitutions most favorably convert one catalytic background into the other. This combination of route diversity and energetic bias is more consistent with a constrained but traversable landscape than with a strictly irreversible transition.

A further caveat concerns the charge model used for the reactive region. In the present parameterization, the resonance charges assigned to sulfur in cysteine and selenium in selenocysteine were effectively identical in the Maestro-derived setup. This result is unlikely to imply electronic equivalence between the two chalcogens. It more plausibly reflects limitations of the charge-derivation protocol and of the underlying force-field treatment for selenium, which remains less extensively parameterized than sulfur. Because sulfur and selenium differ in size, polarizability, and electronegativity, refined quantum-derived charges or experimental calibration could change the quantitative Sec-versus-Cys barrier differences. The qualitative conclusion is less dependent on that detail because, within a single consistent parameterization, the same catalytic substitution has different energetic consequences in the human and mouse scaffolds.

The ProteinMPNN analysis should therefore be interpreted as exploratory rather than as a second main result. It asks which sequences are compatible with the sampled structures, whereas the EVB calculations ask which substitution steps minimize the catalytic barrier. The limited recovery of the node~25 ancestral sequence indicates that structural compatibility alone is insufficient to reconstruct the historical GPX6 sequence under the present design. One potential limitation is the asymmetry in structural starting points, since mouse GPX6 was represented by an experimental crystal structure (PDB 7FC2), whereas human GPX6 relied on an AlphaFold-predicted model. This difference may bias sequence generation toward the mouse background and should be addressed in future work with matched experimental or consistently modeled structures.

Overall, these findings indicate that the mutational pathways linking modern GPX6 homologs to their common ancestor are diverse but not arbitrary. Future studies using deeper variant libraries, explicit functional selection pressures, and protein language models trained or refined with epistatic information may determine whether stronger interactions emerge outside the focused 20-residue library and whether they improve ancestral sequence reconstruction. Such insights would not only advance our understanding of GPX6 evolution but also guide enzyme-engineering strategies that combine sequence models with experimentally or computationally measured epistasis.\cite{Tran2026Science}
